\documentclass[11pt]{article}
\usepackage[margin=25mm]{geometry}
\usepackage{fontspec}
\setmainfont{Latin Modern Roman}
\usepackage{amsmath,amssymb,mathtools}
\usepackage{unicode-math}
\setmathfont{Latin Modern Math}
\usepackage{xurl,hyperref}
\hypersetup{colorlinks=true,urlcolor=blue}
\setcounter{secnumdepth}{0}
\setlength{\parindent}{0pt}
\setlength{\parskip}{5pt}
\setlength{\emergencystretch}{3em}
\newcommand{\sref}[2]{\hyperlink{source:#1:#2}{[S#2]}}
\newcommand{\cataloguescope}[1]{\paragraph{Recorded target.}#1}
\begin{document}
\section{Research Problems in Function Theory — Problem 6.82}\label{statement}

\subsection{Definitions and mathematical statement}

\paragraph{Shared notation.}$\mathbb N=\{1,2,\ldots\}$, and $\mathbb Z,\mathbb Q,\mathbb R,\mathbb C$ denote the integers, rationals, reals and complex numbers. Any other conventions are specified locally.


Write $\mathbb D=\{z\in\mathbb C:|z|<1\}$ and $\mathbb T=\partial\mathbb D$. The class $S$ consists of holomorphic injective functions $f:\mathbb D\to\mathbb C$ normalized by $f(0)=0$ and $f'(0)=1$, so $f(z)=z+\sum_{n=2}^{\infty}a_nz^n$. The inverse in this question is the local inverse germ at $0$, analytically continued to all of $\mathbb D$. It is required to be holomorphic and injective there, but its image need not lie in $\mathbb D$. Let $\sigma$ be the subclass of $S$ whose inverse germs have the continuation just described. Put $a_2^*=\sup_{f\in\sigma}|a_2|$ and $d_n(f)=\bigl||a_{n+1}|-|a_n|\bigr|$.
\paragraph{Statement recorded in the supplied source.}
Is $a_2^*=\sqrt2$? Is there an absolute $A$ with $0<A<1$ such that $d_n(f)\le A$ for every $f\in\sigma$ and every $n\ge2$? Alternatively, determine $\sup_{f\in\sigma}d_n(f)$ for each $n\ge2$. \sref{statement}{2}
Update 6.82 reports verification of the second-coefficient conjecture for certain subclasses only. \sref{statement}{2}
\subsection{Short English statement}
Find the sharp second coefficient for bi-univalent maps, and decide whether all successive coefficient magnitudes have a uniform gap bound strictly below one.
\subsection{Sources}
\begin{description}
\item[\hypertarget{source:statement:1}{S1}] ULAM.AI and the UnsolvedMath Contributors, \emph{UnsolvedMath: A Curated Collection of Open Mathematics Problems}, catalogue entry. CC BY 4.0. \url{https://huggingface.co/datasets/ulamai/UnsolvedMath}.
\item[\hypertarget{source:statement:2}{S2}] W. K. Hayman and E. F. Lingham, \emph{Research Problems in Function Theory} (2018), Chapter 6, Problem 6.82 and its complete update; Chapter 6 initial notation. \url{https://arxiv.org/abs/1809.07200}.
\end{description}
\label{end:statement}
\end{document}
